%% notes-tex/modeling/scaling-laws/scaling-laws.tex
%% [[scaling-laws]]
%% https://github.com/Mjvolk3/torchcell/tree/main/notes-tex/modeling/scaling-laws/scaling-laws.tex
%%
%% Build:  make         -> scaling-laws.pdf
%%         make check   -> formatting / provenance / citation gate
%%         make plots   -> re-convert the panel from ASSET_IMAGES_DIR/041-scaling-laws
%%
%% Figure 1 and Table 1 are written by
%% experiments/041-scaling-laws/scripts/scaling_law_forms.py. Nothing in this document
%% is measured on a torchcell model; Section 5 is a plan and says so.

\documentclass[10pt,a4paper]{article}
\usepackage[draft]{tcdoc}
\usepackage{float}
\usepackage{amsmath}
\usepackage{amssymb}

\emergencystretch=3.5em

\title{\vspace{-1.5em}Scaling Laws for a Genotype-to-Phenotype Model\\
  {\large The two published forms, what counts as $N$, $D$ and $C$ when the
   examples are strains rather than tokens, and the smallest ablation that gives a
   curve worth extrapolating}}
\author{Michael Volk}
\date{\today}

\begin{document}
\maketitle
\thispagestyle{empty}

\begin{abstract}
\noindent
A scaling law is a fitted statement that held-out loss falls as a power of model
size $N$, training data $D$ or training compute $C$, and its value is that a curve
fit on runs one can afford predicts runs one cannot. Two forms carry the
literature. Kaplan et al.\ fit each axis alone, $L = (N_c/N)^{\alpha_N}$ and
$L = (D_c/D)^{\alpha_D}$, straight lines on log-log axes. Hoffmann et al.\ fit the
joint form $L(N, D) = E + A N^{-\alpha} + B D^{-\beta}$ with an irreducible floor
$E$, and minimizing it under the transformer budget $C \approx 6ND$ gives the
compute-optimal allocation $N^* \propto C^{a}$, $D^* \propto C^{b}$ with
$a + b = 1$. One figure draws both forms at their published constants, the plateaus
each axis runs into when the other is fixed, the IsoFLOP curves whose minima trace
the frontier, and a synthetic fitting exercise that holds out the largest runs. For
genotype-to-phenotype data the questions that decide whether the curve means
anything are what to count: $D$ is the number of training examples the optimizer
consumes in the unit the loss averages over, not the dimension of an example;
repeated epochs count for less than fresh data; the evaluation set is fixed and
split by the unit that leaks, a gene or a perturbation, not by record; and $E$ has
a physical reading, the replicate-noise floor, which the released replicate
standard deviations estimate independently of any fit. No torchcell run has been
fit to either form; the last section is the plan for doing so on the trigenic
build, where data, not parameters, is the axis that binds.
\end{abstract}

\vspace{0.5em}
\tccontents
\clearpage

%% notes-tex/modeling/scaling-laws/sections/1-forms.tex

\section{The two forms\secstatus{todo}}
\label{sec:forms}

A \term{scaling law} is a fitted relation between a model's held-out loss and one
or more of the three quantities that a training run spends: the number of trainable
parameters $N$, the amount of training data $D$, and the training compute $C$. The
claim behind the name is that, over the range a family of models was trained on,
the loss falls as a power of each quantity, so that a line fit on log-log axes to
the runs one can afford predicts the loss of a run one cannot. The literature
carries two functional forms, and the difference between them is a floor.

\notesec{Kaplan: one axis at a time}
Kaplan et al.\ (arXiv:2001.08361)\external{arXiv:2001.08361, not in the mirror}
fit each axis separately, with the other axes made large enough not to limit:
\begin{equation}
  L(N) = \left(\frac{N_c}{N}\right)^{\alpha_N}, \qquad
  L(D) = \left(\frac{D_c}{D}\right)^{\alpha_D}.
  \label{eq:kaplan}
\end{equation}
Here $L$ is the \term{test loss}, the average loss on a fixed held-out evaluation
set, cross-entropy per token for a language model. $N_c$ and $D_c$ are fitted
constants in units of parameters and data; they set where each curve sits and have
no meaning of their own. $\alpha_N$ and $\alpha_D$ are the \term{scaling exponents},
the slopes on log-log axes, and the larger one names the axis along which a decade
of spending buys more. For language models Kaplan reports $\alpha_N \approx 0.076$
and $\alpha_D \approx 0.095$. Compute enters through the transformer rule
\begin{equation}
  C \approx 6 N D \ \text{FLOPs},
  \label{eq:sixnd}
\end{equation}
about two floating-point operations per parameter per training token for the
forward pass and four for the backward pass. $N$ in that rule and in
Equation~\ref{eq:kaplan} counts \emph{non-embedding} parameters; including the
embedding matrix bends the curve at small $N$, because the embedding scales with the
vocabulary rather than with the capacity the loss responds to.

\notesec{Chinchilla: a joint form with a floor}
Hoffmann et al.\ (arXiv:2203.15556)\external{arXiv:2203.15556, not in the mirror}
fit both axes at once and let the loss bottom out:
\begin{equation}
  L(N, D) = E + \frac{A}{N^{\alpha}} + \frac{B}{D^{\beta}}.
  \label{eq:chinchilla}
\end{equation}
$E$ is the \term{irreducible loss}, the value approached with infinite parameters
and infinite data; it comes from noise or inherent unpredictability in the data,
the entropy of text for a language model. $A$ and $B$ are the coefficients of the
model-size and data terms, and $\alpha$, $\beta$ are the exponents of the joint fit.
Their Approach~3 fit gives $E = 1.69$, $A = 406.4$, $B = 410.7$, $\alpha = 0.34$,
$\beta = 0.28$ on their training runs. With a floor present, a pure power law is
wrong at the large end: on log-log axes $L$ bends toward $E$ while $L - E$ stays
straight, which is the diagnostic panel~c of Figure~\ref{fig:forms} draws.

\notesec{Compute-optimal allocation}
Minimizing Equation~\ref{eq:chinchilla} subject to the budget $C = 6ND$ gives the
size and data a fixed budget should buy. Writing $D = C/(6N)$ and setting the
derivative in $N$ to zero,
\begin{equation}
  N^*(C) = G \left(\frac{C}{6}\right)^{a}, \qquad
  D^*(C) = \frac{C}{6 N^*(C)}, \qquad
  G = \left(\frac{\alpha A}{\beta B}\right)^{\frac{1}{\alpha + \beta}}, \qquad
  a = \frac{\beta}{\alpha + \beta}, \quad b = \frac{\alpha}{\alpha + \beta},
  \label{eq:optimum}
\end{equation}
so $N^* \propto C^{a}$ and $D^* \propto C^{b}$ with $a + b = 1$. The Hoffmann
constants give $a = 0.452$ and $b = 0.548$, parameters and data growing at nearly
the same rate; Kaplan's earlier conclusion that model size should grow much faster
than data came from learning-rate schedules not matched to each run's length, which
Hoffmann corrected. The \term{compute-optimal frontier} is the curve
$(N^*(C), D^*(C))$ in the $(N, D)$ plane, and an \term{IsoFLOP curve} is the loss
along one budget line $6ND = C$, a U in $N$ whose minimum sits on the frontier.
\src{experiments/041-scaling-laws/scripts/scaling\_law\_forms.py, ChinchillaFit.a / .b / .G}

\notesec{The same runs, a different fit}
Besiroglu et al.\ (arXiv:2404.10102)\external{arXiv:2404.10102, not in the
mirror} refit the Approach~3 data and got $E = 1.82$, $A = 482$, $B = 2085$,
$\alpha = 0.35$, $\beta = 0.37$, which moves the allocation to $a = 0.513$ and
$b = 0.487$. The same runs support two frontiers that differ visibly over the
budgets that matter (panel~h), which is the standing argument for reporting an
interval on every exponent rather than a point.

%% SOURCE: experiments/041-scaling-laws/scripts/scaling_law_forms.py (write_constants_table)
%% Published constants typed from the papers named in the Fit records of that script;
%% a and b are derived from alpha and beta there. Never hand-edit this file.
\begin{table}[htbp]
\centering\footnotesize
\caption{Published constants of the two forms, and the allocation exponents they imply. Kaplan's $N_c$ and $D_c$ set the height of a pure power law and have no meaning on their own; Hoffmann's and Besiroglu's rows are two fits of the same Chinchilla training runs, and the difference between them is the point of panel h. Every value is from outside our holdings.}
\label{tab:constants}
\begin{tabular}{llrl}
\toprule
fit & symbol & value & units \\
\midrule
Kaplan 2020 & $N_c$ & $8.8\times 10^{13}$ & parameters \\
 & $\alpha_N$ & $0.076$ & -- \\
 & $D_c$ & $5.4\times 10^{13}$ & tokens \\
 & $\alpha_D$ & $0.095$ & -- \\
\addlinespace
Hoffmann 2022 & $E$ & $1.69$ & loss \\
 & $A$ & $406.4$ & loss $\cdot$ parameters$^{\alpha}$ \\
 & $B$ & $410.7$ & loss $\cdot$ tokens$^{\beta}$ \\
 & $\alpha$ & $0.34$ & -- \\
 & $\beta$ & $0.28$ & -- \\
 & $a = \beta / (\alpha + \beta)$ & $0.452$ & derived \\
 & $b = \alpha / (\alpha + \beta)$ & $0.548$ & derived \\
\addlinespace
Besiroglu 2024 & $E$ & $1.8172$ & loss \\
 & $A$ & $482.01$ & loss $\cdot$ parameters$^{\alpha}$ \\
 & $B$ & $2085.43$ & loss $\cdot$ tokens$^{\beta}$ \\
 & $\alpha$ & $0.3478$ & -- \\
 & $\beta$ & $0.3658$ & -- \\
 & $a = \beta / (\alpha + \beta)$ & $0.513$ & derived \\
 & $b = \alpha / (\alpha + \beta)$ & $0.487$ & derived \\
\bottomrule
\end{tabular}
\end{table}


\notesec{Symbols}
Three names in the fitting code of Section~\ref{sec:ablation} clash with the
symbols above and are listed here so the clash is visible. Table~\ref{tab:symbols}
is hand-authored.

%% SOURCE: hand-authored; the symbols of eqs. 1-4 and of the estimator in
%% experiments/041-scaling-laws/scripts/scaling_law_forms.py (fit_floored).
\begin{table}[htbp]
\centering\footnotesize
\caption{Every symbol in Equations~\ref{eq:kaplan} to \ref{eq:optimum} and in the
estimator, with the three names that collide.}
\label{tab:symbols}
\begin{tabular}{lp{0.78\textwidth}}
\toprule
symbol & meaning \\
\midrule
$L$ & test loss: average loss on a fixed, separate held-out set \\
$N$ & trainable parameters; Kaplan counts non-embedding parameters only \\
$D$ & training data the optimizer consumed, in the unit the loss averages over (tokens, or records here) \\
$C$ & training compute in floating-point operations \\
$6$ & FLOPs per parameter per training example for a transformer, forward plus backward \\
$N_c$, $D_c$ & Kaplan's fitted scale constants, units of parameters and of data \\
$\alpha_N$, $\alpha_D$ & Kaplan's single-axis exponents \\
$E$ & irreducible loss, the floor \\
$A$, $B$ & coefficients of the finite-size and finite-data terms \\
$\alpha$, $\beta$ & exponents of the joint fit \\
$N^*$, $D^*$ & compute-optimal size and data for a budget $C$ \\
$a$, $b$ & allocation exponents, $N^* \propto C^a$, $D^* \propto C^b$, $a + b = 1$ \\
$G$ & prefactor of the optimum, Equation~\ref{eq:optimum} \\
\addlinespace
\texttt{log\_a}, \texttt{log\_e} & in the estimator: $\log A$ and $\log E$, fit in log space so both stay positive; not the allocation exponents $a$, $b$ \\
\texttt{alpha} & in the estimator: the joint-fit $\alpha$ \\
\texttt{r} & the residual, predicted log-loss minus observed log-loss \\
\texttt{starts} & the grid of initial guesses the optimizer is run from, because the floored fit has several local minima \\
\bottomrule
\end{tabular}
\end{table}

%% notes-tex/modeling/scaling-laws/sections/2-figure.tex

\section{Reading the figure\secstatus{todo}}
\label{sec:figure}

Figure~\ref{fig:forms} draws both forms at the published constants of
Table~\ref{tab:constants}. Every panel is the function, not a measurement; the one
panel with points, panel~i, draws runs from a known curve to show the fitting
recipe on data whose answer is known. The rows go from one axis at a time, to the
joint surface, to what the surface implies for spending.

\tcfig{figures/scaling_law_forms.pdf}{%
The two scaling-law forms and what follows from them.
\textbf{a, b}, Kaplan's single-axis laws, Equation~\ref{eq:kaplan}, with the
published exponent and one smaller and one larger. On log-log axes the exponent
is the slope, and every curve passes through $L = 1$ at $N = N_c$ or $D = D_c$.
\textbf{c}, the irreducible floor: $E + AN^{-\alpha}$ at the Hoffmann constants
bends toward $E$ (dotted) while $L - E$ stays a straight line.
\textbf{d}, the joint form, Equation~\ref{eq:chinchilla}, as a surface over
$\log_{10} N$ and $\log_{10} D$; the amber diagonals are lines of constant compute
$6ND = C$ and the brick line is the compute-optimal frontier, where each diagonal
touches its lowest contour.
\textbf{e}, cuts of the surface at fixed $D$: more parameters stop paying once the
curve reaches the plateau $E + BD^{-\beta}$ (dotted, per curve).
\textbf{f}, cuts at fixed $N$: more data stops paying at $E + AN^{-\alpha}$.
\textbf{g}, IsoFLOP curves, the loss along one diagonal of d, drawn against $N$
with $D = C/(6N)$; each U has a minimum (marked) at $N^*(C)$, and the minima lie
on the frontier (dashed).
\textbf{h}, the allocation, Equation~\ref{eq:optimum}: $N^*$ and $D^*$ against
$C$ on log-log axes, under the Hoffmann fit (solid) and the Besiroglu refit of the
same runs (dashed); the slopes are $a$ and $b$.
\textbf{i}, a synthetic fitting exercise. Fifteen runs, three seeds at each of five
sizes from $10^6$ to $10^{8.67}$, are drawn from $E + AN^{-\alpha}$ at the Hoffmann
constants with 2\% multiplicative noise, and two sizes above are held out. A pure
power law (lilac) and a floored power law (wheat) are fit on the fifteen and
extrapolated, with 5 to 95\% bootstrap bands over runs. The generating curve is
dashed.
}{fig:forms}
\src{experiments/041-scaling-laws/scripts/scaling\_law\_forms.py}

\notesec{What to look for in each row}
In the top row the exponent is everything: panel~a shows that a change of
$\alpha_N$ from 0.076 to 0.15 moves the loss predicted at $N = 10^{13}$ by less
than it moves the loss at $10^5$, since every curve is pinned at $N_c$, but an
exponent fit at small $N$ and used at large $N$ moves the forecast by the full
difference in slope times the decades extrapolated. Panel~c is the diagnostic to
run before trusting any straight-line fit: if the largest runs sit below the line,
the floor has been reached or something else has become the bottleneck, and the
curve to fit is $L - E$.

In the middle row, every cut plateaus, and the height of the plateau is set by the
axis held fixed. Panel~e is the picture of a data-limited project: at
$D = 10^8$ examples the Hoffmann surface flattens at $L = 4.1$ no matter how many
parameters are added, and reaching $L = 3$ needs ten times the data before it
needs a larger model. Panel~f is the mirror image. The dotted plateau lines are
$E + BD^{-\beta}$ and $E + AN^{-\alpha}$, and the gap between a plateau and $E$ is
what the other axis can still buy.

In the bottom row the question is spending. The IsoFLOP minima in panel~g sit at
$N^* = 2.3\times10^8$, $6.4\times10^8$, $1.8\times10^9$ and $5.2\times10^9$
parameters for budgets of $10^{19}$ to $10^{22}$ FLOPs, which is 0.45 decades of
size per decade of compute, the slope $a$ of panel~h. Panel~h also shows the cost
of a fitting choice: the same runs give $a = 0.45$ under one fit and $0.51$ under
another, and at $10^{24}$ FLOPs those two lines are a factor of about two apart
in $N^*$.
\src{experiments/041-scaling-laws/results/scaling\_law\_forms\_summary.json, isoflop\_optima\_N, fits}

Panel~i is the exercise that decides whether a curve is worth extrapolating. The
pure power law fits the fifteen small runs well enough to pass by eye and gives
$\alpha = 0.149$ with a 90\% bootstrap interval of 0.134 to 0.164; it predicts
losses of 1.60 and 1.27 at the two held-out sizes, where the observed means are
1.99 and 1.85, and its band at $10^{10}$ excludes every held-out point. The floored
fit recovers $\alpha = 0.348$ against a generating value of 0.34, with an interval
of 0.316 to 0.376, and $E = 1.71$ against 1.69, and predicts 1.96 and 1.86 at the
held-out sizes. Nothing in the fitted range distinguishes the two models; the
held-out runs do. That is the reason the largest runs are held out rather than
fit.
\src{experiments/041-scaling-laws/results/scaling\_law\_forms\_summary.json, synthetic\_fits}

\notesec{The animated versions}
The Dendron note
\file{notes/experiments.041-scaling-laws.scripts.scaling_law_forms.md} embeds eight
GIFs the same script writes, one sweep each, with the equation, the constants and the
swept value in every frame's title: $D$ rising under the $N$-curve of panel~e, $N$
rising under the $D$-curve of panel~f, $C$ rising under the IsoFLOP curve of panel~g
with the minimum tracing the frontier, $\alpha_N$ moving while the curve is pinned at
the center of the fitted range so that only the forecast moves, $E$ rising under a
fixed $AN^{-\alpha}$ so the bend of panel~c appears, one budget line sliding across
the surface of panel~d, one bootstrap resample and refit per frame accumulating into
the band of panel~i, and $\beta$ moving so the allocation exponents of panel~h move
with it.

%% notes-tex/modeling/scaling-laws/sections/3-counting.tex

\section{Counting $N$, $D$ and $C$ when the examples are strains\secstatus{todo}}
\label{sec:counting}

The forms were fit on language models, where the units are settled: a parameter
is a weight, a datum is a token, and compute is $6ND$. For a model whose examples
are genotype, environment and phenotype records, each unit has to be chosen, and a
curve is only as defensible as the choice. The rule that resolves most of the
cases is that $D$ counts in the unit the loss averages over, and that the unit is
held fixed across every run in the sweep.

\notesec{$D$ is instances consumed, not dimensions}
$D$ is the number of training examples the optimizer has processed when the loss is
measured: the number of records in the training subset times the number of
epochs, in the same unit the loss is averaged over. For a language model that unit
is the token because the cross-entropy is per token; for a model trained on one
loss per record it is the record. The dimension of an example, meaning the number
of features in a record, the number of genes in the genome the model attends over,
or the length of a sequence embedding, is not $D$. Dimension changes what a record
is and how much compute each record costs, and it changes $N$ through the width of
the input layer, but a sweep over dimension is a sweep over architecture, not over
data, and belongs on a different axis of the ablation. Two consequences follow.

\begin{itemize}
\item A record with several supervised labels, a genotype measured on several
phenotypes, counts as one record under a summed loss and as one example per label
under a per-label loss. Either is a valid unit; what breaks the curve is changing
the unit between runs, or reporting a loss averaged over one unit against a $D$
counted in another.
\item A gene-token model that reads a genome as a set of tokens does not consume
$D \times 6{,}000$ tokens in the Kaplan sense: the loss is still one number per
record, and the tokens inside a record are the dimension, not the data. The
per-record compute does scale with that token count, which is why $C$ is measured
rather than computed for such a model (below).
\end{itemize}

\notesec{Repeated epochs count for less}
Language models are trained for about one epoch, so $D$ equals the corpus. A
genotype-to-phenotype model trains for many epochs over a fixed corpus, and a
repeated record is not worth a fresh one. Muennighoff et al.\
(arXiv:2305.16264)\external{arXiv:2305.16264, not in the mirror} fit the
Chinchilla form with an \term{effective data} term in which repeated tokens
decay in value: up to about four epochs the loss is close to what the same number
of fresh tokens would give, and beyond that the return on repetition falls off
quickly. The usable statement for a data-limited project is that $D$ in
Equation~\ref{eq:chinchilla} is the number of distinct records, and the epoch
count is a training hyperparameter held fixed across the sweep, not a multiplier
on $D$. A $D$-sweep is then a sweep over the size of the distinct training set,
which is what asks the question the project needs answered: whether a new
experiment or a bigger model will pay more.

\notesec{Subsample by the unit that leaks}
A $D$-sweep draws nested subsets of the training set, 1/64, 1/16, 1/4 and the
whole, against one fixed evaluation set. If test genotypes share genes or
perturbations with training genotypes, adding data partly measures memorization of
those shared parts, and the apparent exponent is inflated. The subsets are
therefore drawn by the unit that leaks: by gene, or by the query pair a screen was
built around, so that a smaller subset removes whole genes or whole screens rather
than thinning records uniformly. The 025 additive-baselines document showed the
size of this effect on the trigenic build, where the additive ridge falls from
0.400 test Pearson under a random-over-records split to 0.185 under a
query-pair-disjoint one; a $D$-sweep drawn by record on that build would scale
the first number, not the second.

\notesec{$N$: count what the loss responds to}
$N$ is the number of trainable parameters that the loss responds to as capacity.
Kaplan's exclusion of the embedding matrix has a direct analog: a gene-token model
carries an embedding table with one row per gene, whose size is set by the genome
rather than by the architecture, and that table should be reported separately
from the rest, with the curve fit on the rest. Sizes are then produced by scaling
width and depth together in a fixed ratio, five to seven of them spaced evenly in
$\log N$ over at least one and a half to two decades.

\notesec{$C$: measure it}
The $6ND$ rule is transformer-specific and per token. For a model whose per-record
cost depends on how many gene tokens a record carries, on graph message passing,
or on a sequence encoder run once per gene, $C$ is read from a profiler for one
training step at each size and multiplied by the number of steps.
\texttt{torch.utils.flop\_counter.FlopCounterMode} does this for a PyTorch model
without instrumentation. The IsoFLOP construction of panel~g then uses the measured
cost per record in place of $6N$.

\notesec{$E$ has a physical reading}
For a language model $E$ is the entropy of text and is not independently
measurable. For a phenotype measured with replicates it is. Under a mean-squared
loss the best possible predictor has expected error equal to the measurement-noise
variance, so the replicate variance of the evaluation set is a lower bound on $E$
in the loss's own units, and it is available before any model is trained. The
released standard-deviation columns of the Costanzo and Kuzmin screens, whose
replicate structure is recorded in
\file{notes/torchcell.datasets.scerevisiae.costanzo2016.noise-computation.md},
give the per-record $\sigma$ from which that bound is computed. A fitted $E$ that
falls below it says the fit is wrong, and one far above it says the model family
is leaving something on the table that neither more data nor more parameters will
recover. Which of those holds on the trigenic build has not been measured.

%% notes-tex/modeling/scaling-laws/sections/4-ablation.tex

\section{A minimal ablation, and what makes its curve defensible\secstatus{todo}}
\label{sec:ablation}

A Kaplan-style ablation on a modest allocation is three sweeps and four controls.
The sweeps produce the points; the controls decide whether a line through them
means anything.

\notesec{The three sweeps}
\begin{enumerate}
\item \textbf{Parameters.} One architecture family, width and depth scaled in a
fixed ratio, five to seven sizes spaced evenly in $\log N$ over at least one and a
half to two decades, each trained on the full training set until the loss stops
improving. This gives $L(N)$ with data not limiting, the row of panel~a.
\item \textbf{Data.} One model large enough that data is the bottleneck, trained
on nested subsets of the training set at 1/64, 1/16, 1/4 and 1, drawn by the unit
that leaks (Section~\ref{sec:counting}), against one fixed evaluation set. This
gives $L(D)$, the row of panel~b, and for a scientific dataset it is the sweep that
matters most, since $D$ cannot simply be increased.
\item \textbf{Compute.} Three or four budgets $C$. At each, several sizes $N$ with
$D$ set by the budget, $D = C / c_N$ where $c_N$ is the measured cost per record at
that size. Each budget gives a U in $N$, its minimum is $N^*(C)$, and a line
through the minima in log-log is the frontier of panel~h. This is the Chinchilla
IsoFLOP construction, and it is the only one of the three that answers how a fixed
budget should be split.
\end{enumerate}

\notesec{The four controls}
\begin{itemize}
\item \textbf{A learning rate that suits each size.} Kaplan's runs used a cosine
schedule whose length did not match each run's length, and the conclusion about
$N$ versus $D$ changed when Hoffmann matched them. Either tune the learning rate per
size, or use $\mu$P (Yang et al., arXiv:2203.03466)\external{arXiv:2203.03466,
not in the mirror}, under which the optimum transfers across widths so one tuning
at the smallest size serves the sweep.
\item \textbf{Validation loss, not a downstream metric.} Loss is smooth in $N$
and $D$; a correlation or an accuracy jumps and plateaus, and a power law fit to
one is a fit to the plateau's edge.
\item \textbf{Seeds at the small sizes.} Two or three seeds at each of the small
sizes, where they are cheap, give the noise floor the bootstrap needs. Panel~i uses
three.
\item \textbf{Compute counted properly.} $6ND$ for a transformer over tokens;
a profiler for anything else.
\end{itemize}

\notesec{Fitting}
For a pure power law, a linear regression of $\log L$ on $\log N$ is the fit, and
the slope is the exponent. With a floor the fit is nonlinear least squares in log
space over $(\log A, \log E, \alpha)$, run from a grid of starting points and
keeping the best, because the objective has several local minima; fitting the
logarithms keeps $A$ and $E$ positive without constraints. The estimator in
\file{experiments/041-scaling-laws/scripts/scaling_law_forms.py} (function
\texttt{fit\_floored}) does this with 27 starts and a squared residual. Hoffmann's
appendix uses a Huber residual, which stops a few bad runs from steering the
exponent; it is the right choice once real runs are being fit and is a one-line
change there. The joint form adds $(\log B, \beta)$ to the parameter vector and is
otherwise the same fit.

\notesec{What makes the curve defensible}
\begin{itemize}
\item \textbf{Hold out the largest runs.} Fit on the smaller runs and check whether
the fit predicts the largest. Panel~i is the case where this is the only test that
separates a right model from a wrong one.
\item \textbf{Bootstrap over runs.} Resample runs with replacement, refit, and
report an interval on every exponent. A point estimate of $\alpha$ is not a
result, and panel~h shows two published point estimates from one set of runs.
\item \textbf{Look for curvature.} A bend at the large end means $E$ is near, or
that hyperparameters or data have become the bottleneck. Either way the straight
line stops there.
\item \textbf{Check $E$ against the replicate floor.} When the loss is a squared
error on a replicated measurement, the replicate variance is an estimate of $E$
that owes nothing to the fit (Section~\ref{sec:counting}).
\end{itemize}

%% notes-tex/modeling/scaling-laws/sections/5-torchcell.tex

\section{What an ablation looks like on the trigenic build\secstatus{todo}}
\label{sec:torchcell}

\begin{keybox}
\textbf{Measured.} No number here is a torchcell result. The constants
are the published ones and panel~i is synthetic.

\textbf{Not measured.} Every number a scaling section of the manuscript or thesis
would state: the exponents, the floor, and whether a curve fit on small runs
predicts a large one. What follows is the design, with the choices each sweep
forces written down so a run can be launched against it.
\end{keybox}

\notesec{The axes available}
The trigenic build of experiment 025 holds 376{,}732 records with bit-identical
labels to the 010 build, organized as 420 recurring query pairs plus one array
gene each, and it already carries the two splits an ablation needs: the published
random-over-records split and a query-pair-disjoint split under which no query
double is shared between parts. The equivariant cell graph transformer trained on
it is one architecture family whose width and depth can be scaled together. The
axes are therefore available now; what is missing is the sweep.

\notesec{The data sweep first}
Data is the binding axis here, so the $D$-sweep is the one to run before any
other. Nested subsets of the disjoint-split training set at 1/64, 1/16, 1/4 and 1,
drawn by query pair so that a smaller subset removes whole screens, against the
one fixed disjoint test set; a model at the current configuration's size; the
loss recorded per epoch and the validation-selected checkpoint scored on test.
Four points give a slope and, with the held-out largest subset, the first
extrapolation test. The reading the sweep is for: whether the loss is still
falling with $D$ at the full set, which says a new screen is worth more than a
larger model, or has flattened, which says the opposite.

\notesec{The parameter sweep}
Five sizes of the transformer, width and depth in the fixed ratio of the current
configuration, spaced evenly in $\log N$ over about two decades, with the gene
embedding table counted separately and the curve fit on the rest; three seeds at
the two smallest sizes; learning rate tuned per size, or transferred under $\mu$P
once the parametrization is checked against the model code. Every run on the full
disjoint training set, so that this sweep reads $L(N)$ with data as fixed as it
can be made.

\notesec{The floor from the replicates}
Before either sweep, the replicate standard deviations of the test records give
the noise-variance bound on $E$ in the loss's units. Whether the trigenic
interaction score's released uncertainty is a sample standard deviation over
colonies or a bootstrap standard error, and over how many colonies, is already
recorded for the Kuzmin and Costanzo screens in the noise-computation note, and
the bound follows from it. That number is the first thing to set beside a fitted
$E$.

\notesec{Compute}
The cost per record of the transformer depends on how many gene tokens a record
carries and on the graph penalty, so $C$ is read from a profiler at each size
rather than from $6ND$. The IsoFLOP sweep is third in priority: it answers how to
split a budget, which matters once the first two sweeps have shown that the
budget buys anything.

\notesec{What the section in the manuscript would then say}
A scaling section is one figure, three panels, $L$ against $D$, against $N$ and
against $C$, each with its fitted exponent and interval, the held-out largest run
marked, and $E$ drawn against the replicate floor. Until the sweeps have run, that
figure is a plan, and the sweeps above are its specification.

%% notes-tex/modeling/scaling-laws/sections/6-expression.tex
%% SOURCE: tables/floors.tex, tables/budget_curve.tex and the macros of tables/expression_numbers.tex
%% are generated by experiments/041-scaling-laws/scripts/expression_record.py from the committed 019
%% result files it names (expression_ceiling_all, morphology_noise_ceiling, short_budget_spread,
%% loss_min_vs_pearson_peak, joint_checkpoint_readout, baselines_both_label, v21_readout and
%% prediction_shrinkage_probe). The parameter-count reading and the two-store comparison are quoted
%% from notes-tex/019-simb-multimodal-expression/sections/0-strand.tex and
%% notes-tex/figure-3-gate/ (branch multimodal-phenotype-retrospective) and are named as such.
%% GENERATED FILE -- do not hand-edit.
%% SOURCE: experiments/041-scaling-laws/scripts/expression_record.py, reading experiments/019-simb-multimodal/results/the files named in the script docstring
\newcommand{\budgetFirstEpochs}{350}
\newcommand{\budgetFirstMean}{0.120}
\newcommand{\budgetLastEpochs}{9,900}
\newcommand{\budgetLastMean}{0.197}
\newcommand{\budgetThousandMean}{0.161}
\newcommand{\budgetGainPastThousand}{0.036}
\newcommand{\lossLiveRuns}{27}
\newcommand{\lossMinEpochMedian}{481}
\newcommand{\lossMinEpochMax}{3,097}
\newcommand{\lossPeakEpochMedian}{2,433}
\newcommand{\lossFinalExcessPct}{9.0}
\newcommand{\lossAtPeakExcessPct}{6.8}
\newcommand{\lossGainAfterMin}{+0.057}
\newcommand{\jointExprDiff}{-0.0012}
\newcommand{\jointExprPos}{5}
\newcommand{\jointExprN}{11}
\newcommand{\jointExprRef}{0.094}
\newcommand{\jointProtDiff}{-0.0172}
\newcommand{\jointProtPos}{1}
\newcommand{\jointProtMargin}{0.015}
\newcommand{\exprBar}{0.110}
\newcommand{\exprBarSd}{0.017}
\newcommand{\exprBarSeeds}{12}
\newcommand{\vBestArm}{S\_mask}
\newcommand{\vBestMean}{0.094}
\newcommand{\vBestSd}{0.028}
\newcommand{\vRefMean}{0.063}
\newcommand{\vRefSd}{0.037}
\newcommand{\shrinkSpreadMin}{1.4}
\newcommand{\shrinkSpreadMax}{2.2}
\newcommand{\shrinkRuns}{3}
\newcommand{\shrinkProtSpreadMin}{0.5}
\newcommand{\shrinkProtSpreadMax}{0.7}
\newcommand{\shrinkProtRuns}{3}
\newcommand{\morphRealized}{13}
\newcommand{\morphObservedBest}{0.082}


\section{What the expression and proteome rounds already pin down\secstatus{todo}}
\label{sec:expression}

\begin{keybox}
\textbf{Measured.} The replicate floor of every 019 panel in Pearson units; a
budget curve over eight arms from 350 to 9{,}900 epochs; the disagreement between
the validation loss and the metric in \lossLiveRuns\ long runs; a joint-training
null on eleven partitions; the sequence-baseline bar and the best arm of the
twelve-seed round.

\textbf{Not measured.} A parameter sweep and a data sweep on the expression
store. Every 019 round varied one architectural choice at a fixed size on a fixed
store, and the sizes that differ between rounds differ with the store and the
input stack, so no exponent can be read from the record.
\end{keybox}

Experiment 019 trained the equivariant cell graph transformer on the single-deletion
transcriptome (Kemmeren 2014), the single-deletion proteome (Messner 2023) and
morphology (Ohya 2005), alone, jointly and conditioned on one another, across about
twenty rounds. Read against Sections~\ref{sec:counting} and \ref{sec:ablation}, that
record settles the floor, exhibits the compute axis on the wrong statistic, and
leaves both scaling axes unmeasured. Each piece is listed with what it is good for.

\notesec{The floor, per panel}
Table~\ref{tab:floors} is the measurement Section~\ref{sec:counting} asks for before
any sweep, done for four panels. Kemmeren's two independent cultures per deletion
give a test-retest ceiling of a mean Pearson per feature of 0.775 across strains;
Caudal's 29 re-cultured isolates 0.604; Messner's duplicated-origin strains 0.405
against 0.696 from its wild-type replicates, two estimators that disagree because 146
pairs determine neither well; morphology 0.611 from 122 wild-type replicates. These are
in Pearson, the campaign's metric. The bound on $E$ in loss units needs the per-gene
replicate variance, which Kemmeren's store carries as a standard error and a replicate
count per value and which has not been converted; for the other panels only the
pooled reliability is on file.

%% GENERATED FILE -- do not hand-edit.
%% SOURCE: experiments/041-scaling-laws/scripts/expression_record.py, reading experiments/019-simb-multimodal/results/expression_ceiling_all.json and morphology_noise_ceiling.json
\begin{table}[htbp]
\centering\footnotesize
\caption{The replicate ceiling of each 019 panel, the bound on $E$ in Pearson units. Reliability is the per-feature fraction of across-strain variance that is signal (test-retest: the correlation of two independent measurements across strains; decomposition: one minus replicate variance over total variance), averaged over features; the ceiling is the mean of its square root, the best across-strain Pearson a perfect predictor of the signal can reach against a single noisy measurement. Two estimators are given where both exist because they disagree where few pairs determine either.}
\label{tab:floors}
\begin{tabular}{llrr}
\toprule
panel & estimator & reliability & ceiling, mean Pearson per feature \\
\midrule
Kemmeren 2014, expression & test-retest, 82 paired deletions & 0.611 & 0.775 \\
 & decomposition & -- & 0.862 \\
Caudal 2024, expression & test-retest, 29 re-cultured isolates & 0.405 & 0.604 \\
 & decomposition & 0.235 & 0.459 \\
Messner 2023, proteome & test-retest, 146 duplicated-origin ORFs & 0.200 & 0.405 \\
 & decomposition, 388 wild-type replicates & 0.501 & 0.696 \\
Ohya 2005, morphology & decomposition, 122 wild-type replicates & 0.444 & 0.611 \\
\bottomrule
\end{tabular}
\end{table}


\notesec{The compute axis, in Pearson}
Table~\ref{tab:budget-curve} is a budget sweep of the kind panel~c of the figure
draws, read off the training curves of eight arms rather than run as one: the same
runs rescored on the prefix of their own history at ten budgets. The mean rises from
\budgetFirstMean\ at \budgetFirstEpochs\ epochs to \budgetLastMean\ at
\budgetLastEpochs, and flattens: the tenfold budget past 1{,}000 epochs buys
\budgetGainPastThousand. That is the shape a power law with a floor has, and it
cannot be fit as one, because the statistic is a correlation and not a loss. The
document's second control (Section~\ref{sec:ablation}) exists for exactly this case.

%% GENERATED FILE -- do not hand-edit.
%% SOURCE: experiments/041-scaling-laws/scripts/expression_record.py, reading experiments/019-simb-multimodal/results/short_budget_spread.json (W&B torchcell_019_expr_v9)
\begin{table}[htbp]
\centering\footnotesize
\caption{The compute axis as the 019 campaign measured it: the same eight long-budget arms of the v9 round rescored on the prefix of their own curves at each budget (the maximum of a centered five-epoch rolling mean of the validation Pearson per feature over epochs up to the budget). Mean and sd over the eight arms; the spread is an arm spread plus nondeterminism, so it bounds the replicate spread from above. The statistic is Pearson, not the loss, for the reason given in the text.}
\label{tab:budget-curve}
\begin{tabular}{rrrrrr}
\toprule
epochs & arms & mean Pearson per feature & sd & min & max \\
\midrule
350 & 8 & 0.1198 & 0.0096 & 0.1029 & 0.1317 \\
500 & 8 & 0.1223 & 0.0058 & 0.1130 & 0.1317 \\
700 & 8 & 0.1411 & 0.0166 & 0.1207 & 0.1606 \\
1,000 & 8 & 0.1609 & 0.0099 & 0.1485 & 0.1780 \\
1,500 & 8 & 0.1670 & 0.0117 & 0.1536 & 0.1839 \\
2,000 & 8 & 0.1745 & 0.0151 & 0.1547 & 0.1950 \\
2,800 & 8 & 0.1821 & 0.0175 & 0.1629 & 0.2085 \\
4,000 & 8 & 0.1883 & 0.0171 & 0.1663 & 0.2145 \\
6,000 & 8 & 0.1917 & 0.0177 & 0.1663 & 0.2222 \\
9,900 & 8 & 0.1965 & 0.0222 & 0.1663 & 0.2382 \\
\bottomrule
\end{tabular}
\end{table}


\notesec{Why the loss cannot be the objective as it stands}
In the \lossLiveRuns\ long runs of the same round whose loss is not the metric, the
validation loss bottoms out at a median epoch \lossMinEpochMedian\ (the latest at
\lossMinEpochMax) and the Pearson peaks at a median epoch \lossPeakEpochMedian, with
the loss already \lossAtPeakExcessPct\,\% above its minimum and
\lossGainAfterMin\ of Pearson gained between the two; the final loss sits
\lossFinalExcessPct\,\% above its minimum. The shrinkage probe on saved validation
predictions says why: the expression predictions are \shrinkSpreadMin\ to
\shrinkSpreadMax\ times too spread for their correlation on \shrinkRuns\ checkpoints
(the proteome's run the other way, \shrinkProtSpreadMin\ to \shrinkProtSpreadMax\ on
\shrinkProtRuns), and one scalar per run brings the squared error below the per-gene
mean. The loss is
a scale problem riding on a correlation that keeps improving. A scaling fit on the
raw loss would therefore read a calibration artifact as the curve. Two repairs are
available and neither has been run as a sweep: fit on a calibrated loss (the squared
error after the one-scalar shrinkage, which is the probe's statistic), or fit on the
per-gene Pearson with the understanding that it is bounded and will bend near the
floor for a reason that is not $E$.

\notesec{The bar a curve has to start above}
On the store the twelve-seed round trains on, the nearest-neighbor average over
ProtT5 embeddings reaches \exprBar\ $\pm$ \exprBarSd\ validation Pearson per feature
over \exprBarSeeds\ split seeds. The best arm of that round, \texttt{\vBestArm},
reaches \vBestMean\ $\pm$ \vBestSd\ over twelve seeds, and the reference
\vRefMean\ $\pm$ \vRefSd. Morphology realizes \morphRealized\,\% of its ceiling. A
curve fit to a model that sits at or below a trivial predictor measures the predictor,
so on this store the first point of either sweep has to clear the neighbor average
before the second point means anything. On the larger expression store the v13
reference sits within 0.01 of the same baseline (Figure~3 gate document, branch
\file{multimodal-phenotype-retrospective}).

\notesec{The multimodal data axis}
Adding the proteome's strains and labels to the trunk is a move along $D$ on a second
modality, and the joint round measured it on eleven partitions: the joint arm against
the expression-only arm \jointExprDiff\ (\jointExprPos\ of \jointExprN\ partitions
above zero, reference \jointExprRef), and against the proteome-only arm
\jointProtDiff\ (\jointProtPos\ of \jointExprN\ above zero, against a
non-inferiority margin of \jointProtMargin). Pooling a second modality did not move
the curve. Conditioning did, by $+0.03$ to $+0.05$ over the unconditioned arm
(Figure~3 gate document), and stayed below a ridge from the measured modality, which
is the linear baseline for that axis. This is the one data-axis measurement the record
holds, and it says the second modality is worth more as an input than as more
training signal for a shared trunk.

\notesec{What the record says about $N$, and why it is not a curve}
The expression strand's own summary reads 145 runs from 0.12M to 4.7M parameters
with a correlation of $-0.129$ between $\log_{10} N$ and score
(\file{notes-tex/019-simb-multimodal-expression/sections/0-strand.tex}), and the
Figure~3 gate found that no committed script reproduces that subset while all 1{,}182
scored runs give $-0.001$. Either way the runs differ in budget, input stack, head and
objective at once, so the reading is a confounded null, not $L(N)$. The one deliberate
width increase on the fast code, four layers at width 180 against six at width 90,
collapsed to the per-gene-mean predictor on 2 of 2 seeds, which is the first control
of Section~\ref{sec:ablation} failing before the sweep could start: the learning rate
that suits width 90 does not suit width 180.

\notesec{What the record does not say about $D$}
The best per-gene scores of the campaign, the v13 reference at a window mean of 0.20
and 0.17 on split seeds 0 and 1 (\file{v13\_split\_readout.json}), were trained on
about the same number of strains as the twelve-seed round (the fig3\_core expression
rows against the 1{,}103 both-label rows), with the full four-embedding input stack
and a 6{,}000-epoch budget. The gap to the small trunk's 0.06 to 0.10 is therefore
budget and input, not $D$, and the record holds no two points that differ in $D$
alone.

\notesec{The two sweeps, launched on cabbi 2026-10-09}
Both run on the fast recipe with the mask schedule on (the one twelve-seed arm with
zero collapses) at batch 128 and the unscaled rate, one run per RTX6000 card, three
split seeds, W\&B project \file{torchcell\_019\_scaling}, stages \file{scale\_d} and
\file{scale\_n} of \file{igb\_expr\_wave5.slurm} (IGB jobs 2427637 and 2427638). The
statistic is still the per-gene Pearson at the registered window, with the loss curve
logged for the calibrated fit later.

\begin{itemize}[leftmargin=1.2em,itemsep=1pt,topsep=2pt]
\item \textbf{Data.} Nested prefixes of one split-seeded permutation of the 1{,}103
both-label training rows at 1, 1/2, 1/4, 1/8 and 1/16, validation and test untouched
(\file{train\_cgt\_multitask.py}, \texttt{train\_fraction}). The budget is 1{,}200
epochs for every fraction, so the small fractions take fewer optimizer steps (one per
epoch at 1/16); the per-epoch validation curve is what the readout reads and the step
mismatch is part of the result, not hidden by it.
\item \textbf{Parameters.} Widths 45, 90, 180 and 360 at six layers (nine heads
divide all four), the learning rate scaled as $1/\text{width}$ from $3\times10^{-4}$
at 90. Hypothesis (untested): that rule keeps the wider trunks out of the collapse the
width-180 arm fell into at the width-90 rate; it is not $\mu$P, and a width whose
rate is wrong will show as a collapse rather than a point on the curve. Width 360 is
last in the array so an out-of-memory there costs the sweep nothing else.
\end{itemize}

%% notes-tex/modeling/scaling-laws/sections/7-sources.tex

\section{Sources\secstatus{todo}}
\label{sec:sources}

None of the five works below is in the mirror, so each is named by arXiv
identifier and flagged where its numbers are used. Adding them to the library is
a curation decision that has not been made; once it is, the citations here become
\texttt{\textbackslash cite} keys and the flags come off.

\begin{itemize}
\item Kaplan et al., 2020, \emph{Scaling Laws for Neural Language Models},
arXiv:2001.08361. The single-axis forms of Equation~\ref{eq:kaplan}, the
non-embedding parameter count, and the $6ND$ rule.
\item Hoffmann et al., 2022, \emph{Training Compute-Optimal Large Language
Models}, arXiv:2203.15556. The joint form with a floor, the three fitting
approaches (its appendix covers all three, including the Huber residual), the
IsoFLOP construction, and the Approach~3 constants.
\item Besiroglu et al., 2024, \emph{Chinchilla Scaling: A Replication Attempt},
arXiv:2404.10102. The refit of the Approach~3 data and the second set of
constants in Table~\ref{tab:constants}.
\item Yang et al., 2022, \emph{Tensor Programs V: Tuning Large Neural Networks
via Zero-Shot Hyperparameter Transfer}, arXiv:2203.03466. $\mu$P, the
parametrization under which the learning-rate optimum transfers across widths.
\item Muennighoff et al., 2023, \emph{Scaling Data-Constrained Language Models},
arXiv:2305.16264. The effective-data form for repeated epochs.
\end{itemize}

Within the repository: the 025 additive-baselines document
(\file{notes-tex/multimodal/025-additive-baselines/}) for the record count, the
splits and the leakage measurement quoted in Section~\ref{sec:counting}, and
\file{notes/torchcell.datasets.scerevisiae.costanzo2016.noise-computation.md} for
the replicate structure behind the floor estimate.


\end{document}
